% ===== §7 The Neural Plasticity of Perception =====
\section{The Neural Plasticity of Perception}
\label{p23:sec:plasticity}

\noindent
The philosophy of the body established that perceiving is an embodied skill and therefore, in principle, the kind of capacity an education could reach. This section supplies the mechanism. It is the most empirical chapter of the paper, and it carries the second spine's weight, because a claim that perception can be cultivated is idle unless the perceiving system is in fact modifiable by experience in the relevant way. The neuroscience of the last decades has established that it is, and established it at the level that matters: not merely that preferences shift or that judgements are learned, but that the perceptual system itself, down to the tuning of neurons in early sensory cortex, is reorganised by training. The section sets out this evidence, connects it to the predictive account already introduced, and then meets the objection that the neuroscience of the aesthetic measures pleasure rather than perception, an objection this paper must answer with particular care because answering it is where its own use of neuroscience becomes distinctive. It closes by drawing, from the very plasticity it has established, the line between the cultivation of perception and its commodified training, which is the second spine's ethical bound.

\subsection{Plasticity as the empirical premise of cultivation}
\label{p23:sub:plas-premise}

The enabling fact is that neural plasticity is not an exceptional or developmental state of the brain but its ongoing condition. It was once assumed that the perceptual system fixed early, that after a critical period in development the neural mechanisms of early sensory processing were no longer modifiable, and that what improvement adults showed with practice was a matter of higher cognition, of decision and judgement laid over an unchanging perceptual base. This assumption has not survived the evidence. Cortical representations of sensory input remain plastic in the mature brain, and the reorganisation of neural circuits by experience is the normal ongoing state of the nervous system across the lifespan \parencite{gilbert2001}. This is the premise on which an education of perception rests. If the perceptual system were fixed after development, the cultivation of perception could reach only the judgements and associations built on top of a standard perceiving, and the paper's second spine would collapse into a claim about taste rather than sight. Because the perceptual system is plastic throughout life, the cultivation of perception can be what the paper claims it is: a training of perceiving itself.

\subsection{Perceptual learning and cortical reorganisation}
\label{p23:sub:plas-perceptuallearning}

The direct evidence is the phenomenon of perceptual learning: the sustained and robust improvement, through repeated practice or exposure, in the ability to discriminate and identify perceptual features, an improvement that develops gradually and largely without conscious effort and is maintained over months and years \parencite{watanabe2015}. Perceptual learning is not the acquisition of a belief about what one sees; it is an enhancement of the seeing. Its behavioural signature, improved discrimination at the trained task, is accompanied by measurable change in the nervous system at multiple levels, including the tuning of neurons in early sensory cortex, the correlational structure of sensory populations, and the weighting by which sensory evidence is read out for perception \parencite{gilbert2001,watanabe2015}. What training alters is the perceptual apparatus itself, the machinery by which a stimulus becomes a discriminated feature, and not merely a downstream decision about an already-fixed percept. For the argument of this paper the significance is exact. When it is claimed that a person may be brought to perceive the everyday more finely, to discriminate qualities of light, texture, sound, and taste that were previously undifferentiated, the claim names a real and studied form of neural change and not a metaphor of enrichment. The refinement of perception is the perceptual learning of the everyday, and it has the same standing as the perceptual learning studied in the laboratory.

\subsection{Expertise and the reorganised perceiver}
\label{p23:sub:plas-expertise}

The most vivid form of the evidence is the perceptual expert, and the expert cases matter to this paper because they show the endpoint toward which cultivation tends. Sustained, deliberate engagement with a domain reorganises the brain, structurally and functionally, in ways that imaging has documented across musicians, athletes, and other trained perceivers \parencite{watanabe2015}. The trained musician hears a texture of sound where the untrained hears an undifferentiated chord, and the difference is not merely that the musician has words for what they hear; their auditory processing has been reshaped by training so that the discrimination is perceptual. The taster discriminates in a mouthful of wine what the novice cannot detect at all; the trained eye of the radiologist sees, in a field the layperson finds uniform, the salient and the anomalous. These are not additions of knowledge to a common perceptual base. They are reorganisations of perceiving, acquired through a history of engagement, by which the expert comes to inhabit a perceptually richer world than the novice occupies before the same stimulus. The expert perceiver is the demonstration, at the limit, of what the cultivation of everyday perception aims at in small: a person whose sustained attention to a domain has remade their capacity to perceive within it, so that the world answers their trained perception with a richness it withholds from the untrained. What the sommelier is to wine, a cultivated everyday perceiver is to the light of ordinary rooms and the textures of ordinary days.

\subsection{The predictive implementation}
\label{p23:sub:plas-predictive}

These findings are not a loose collection; they are unified by the predictive account of the brain already introduced, and the unification matters because it ties the plasticity of perception to the generativity established in the first spine. If perception is the generation of a predicted world corrected by sensory error \parencite{friston2010}, then perceptual learning is the refinement of the generative model under the pressure of that error, and the reorganisation documented in sensory cortex is the physical form of a model that has been retuned by experience to predict a domain more finely. The expert's richer perceptual world is the world afforded by a more finely differentiated generative model, one whose predictions carve the domain at joints the novice's model does not represent. This connects the second spine to the first at the level of mechanism. The generativity of perception and the plasticity of perception are one fact seen twice: perception generates its world under a model, and cultivation is the retuning of that model, so that to cultivate perception is to reshape the very generative activity by which perception constructs what is perceived. The everyday, once more, is where the retuning most matters, because it is the domain in which an unrefined generative model most readily predicts the familiar away into invisibility, and in which a cultivated model most transforms what is available to be perceived.

\subsection{The objection from pleasure}
\label{p23:sub:plas-neuroaesthetics}

An objection arises precisely from the existing neuroscience of the aesthetic, and it must be met with care, because meeting it is where this paper's use of neuroscience declares its difference. Neuroaesthetics, the field named by Zeki and developed through the aesthetic-triad model of Chatterjee and Vartanian, studies the neural bases of aesthetic experience by way of three interacting systems, the sensory-motor, the emotion-valuation, and the meaning-knowledge, and its characteristic finding locates the experience of beauty in the pleasure centres and reward circuits of the brain, with pleasure treated as a necessary if not sufficient condition of the beautiful \parencite{zeki1999,chatterjee2014}. The objection this invites against the present paper runs as follows: if the neuroscience of the aesthetic shows that aesthetic experience is a matter of reward and valuation, then what any training of the aesthetic trains is preference, liking, the hedonic response to a stimulus, and not perception at all; the talk of cultivating perception collapses into the conditioning of pleasure. The paper's answer turns the structure of the aesthetic triad against the objection rather than denying the triad. It grants that the emotion-valuation system is engaged in aesthetic experience and that reward is a real component of it. But the reduction of the aesthetic to reward conflates distinct systems that the triad itself separates, and it conflates, within the valuation system, the wanting and the liking that a careful reward neuroscience has shown to be dissociable \parencite{berridge2015,kringelbach2009}. The cultivation this paper is concerned with does not principally engage the emotion-valuation system at all. It engages the sensory-motor system, the perceptual apparatus whose plasticity the preceding sections documented; its object is not the intensification of the pleasure taken in a stimulus but the reorganisation of the perceiving by which the stimulus is discriminated. This is the paper's distinctive claim about the neuroscience of the aesthetic, and it is where a background in the dynamics of perceptual and neural systems earns its place. Where neuroaesthetics has largely measured the reward correlates of preference, asking which stimuli the brain likes, this paper points the neuroscience at a different question, asking how the perceiving itself is reorganised, and locating the aesthetic cultivation of the everyday in the plasticity of perception rather than in the conditioning of reward. The sommelier is not someone trained to like wine more; the sommelier is someone whose perception of wine has been remade. The objection from pleasure holds only against a project that sought to train pleasure, and that is not this project.

\subsection{The line between refinement and commodified training}
\label{p23:sub:plas-redline}

The plasticity that makes cultivation possible makes its abuse possible in the same measure, and the second spine must therefore draw a line within the very trainability it has established. Perception can be reorganised, and it can be reorganised toward the refinement this paper commends or toward an end this paper must mark as its corruption. The refinement of perception enlarges the perceiver's world, deepening the discriminations available to them and returning them to the everyday with a richer capacity to generate its perception. The commodified training of perception does the contrary under the same mechanism: it reorganises perceiving so as to bind it to consumption, training the eye to discriminate the marks of status and the triggers of purchase, tuning the perceptual system to the demands of an attention economy that profits from what it teaches the senses to want. Both are perceptual learning; both retune the generative model; both are, at the level of mechanism, the same neural plasticity. The line between them is not drawn in the neuroscience, which is neutral between them, but in the value theory the paper has already supplied. Refinement is the cultivation of a perception whose value returns to the perceiver and to the shared attending that generated it, a perception drawn toward the register of the real and the good circulation of everyday aesthetic value. Commodified training is the cultivation of a perception whose value is extracted, a perceiving reorganised to serve a circuit of consumption in which the perceiver's own trained senses are turned into an instrument of their capture. The plasticity is one; the two uses of it are told apart by whether the reorganised perception generates value that returns to the perceiver or is harnessed to extract it. This is the ethical bound of the second spine, and it hands the argument forward to the theory of justice, where the social distribution of perceptual cultivation, and the question of whose senses are refined and whose are trained for capture, becomes a question of justice proper (\xref{p23:sec:distribution}).
